% Standalone problem document, generated from the adjacent README.md.
% Compile with XeLaTeX. Mathematical content is maintained in Markdown.
\documentclass[11pt,a4paper]{article}
\usepackage[fontsize=10.5pt]{scrextend}
\usepackage[margin=22mm,headheight=15pt,headsep=9mm,footskip=11mm]{geometry}
\usepackage{fontspec}
\setmainfont{texgyrepagella}[Extension=.otf,UprightFont=*-regular,BoldFont=*-bold,ItalicFont=*-italic,BoldItalicFont=*-bolditalic]
\setsansfont{texgyreheros}[Extension=.otf,UprightFont=*-regular,BoldFont=*-bold,ItalicFont=*-italic,BoldItalicFont=*-bolditalic]
\setmonofont{texgyrecursor}[Extension=.otf,UprightFont=*-regular,BoldFont=*-bold,ItalicFont=*-italic,BoldItalicFont=*-bolditalic,Scale=MatchLowercase]
\usepackage{amsmath,amssymb}
\usepackage{unicode-math}
\setmathfont{texgyrepagella-math.otf}
\usepackage{xcolor,microtype,enumitem,needspace,fancyhdr}
\usepackage{bookmark}
\definecolor{ink}{HTML}{17344B}
\definecolor{accent}{HTML}{197D80}
\definecolor{muted}{HTML}{596773}
\hypersetup{colorlinks=true,linkcolor=accent,urlcolor=accent,pdftitle={MF-07 - Uniform
polynomial bounds for products at joint spectral radius
one},pdfauthor={Open Problems in Numerical Linear Algebra}}
\urlstyle{same}
\setlength{\parindent}{0pt}
\setlength{\parskip}{5pt plus 1pt minus 1pt}
\linespread{1.04}
\setlength{\emergencystretch}{3em}
\widowpenalty=10000
\clubpenalty=10000
\displaywidowpenalty=10000
\allowdisplaybreaks[1]
\setlist{nosep,leftmargin=1.5em}
\providecommand{\tightlist}{\setlength{\itemsep}{0pt}\setlength{\parskip}{0pt}}
\providecommand{\pandocbounded}[1]{#1}
\pagestyle{fancy}
\fancyhf{}
\fancyhead[L]{\sffamily\footnotesize\color{muted}OPEN PROBLEMS IN NUMERICAL LINEAR ALGEBRA}
\fancyhead[R]{\sffamily\bfseries\footnotesize\color{accent}MF-07}
\fancyfoot[L]{\parbox[b]{0.9\textwidth}{\sffamily\fontsize{8}{10}\selectfont\color{muted}Editorial ratings; a bounded search cannot certify that a problem remains unsolved.\\Literature check: 2026-09-11}}
\fancyfoot[R]{\sffamily\footnotesize\color{muted}\thepage}
\renewcommand{\headrulewidth}{0.3pt}
\renewcommand{\headrule}{\hbox to\headwidth{\color{accent}\leaders\hrule height \headrulewidth\hfill}}
\makeatletter
\renewcommand{\section}{\Needspace{5\baselineskip}\@startsection{section}{1}{0pt}{12pt plus 2pt minus 2pt}{4pt}{\sffamily\bfseries\large\color{ink}}}
\renewcommand{\subsection}{\Needspace{4\baselineskip}\@startsection{subsection}{2}{0pt}{9pt plus 1pt minus 1pt}{3pt}{\sffamily\bfseries\normalsize\color{ink}}}
\makeatother
\setcounter{secnumdepth}{0}
\begin{document}
{\sffamily\small\color{muted}Matrix functions and stability\par}
\vspace{3pt}
{\raggedright\hyphenpenalty=10000\exhyphenpenalty=10000\sffamily\bfseries\fontsize{21}{24}\selectfont\color{ink}Uniform
polynomial bounds for products at joint spectral radius one\par}
\vspace{7pt}
{\sffamily\small\textbf{Difficulty:} challenging\quad\textbf{Importance:} interesting
to the community\par}
{\sffamily\small\bfseries\color{accent}Status: Solved\par}
\vspace{4pt}
{\color{accent}\hrule height 0.8pt}
\vspace{6pt}
\textbf{Rating rationale:} Challenging because the constant must be
uniform over all families of a given dimension; community impact spans
transient growth and stability analysis.

\subsection{Resolution --- 2026-09-11}\label{resolution-2026-09-11}

\textbf{Affirmative resolution.} Matthew J. Colbrook's
\href{https://github.com/ajt60gaibb/OpenProblemsInNLA/blob/main/references/colbrook-jsr-growth-2026-09-11/manuscripts/uniform_growth_and_holder.pdf}{complete
manuscript, Theorem 1 and Proposition 5} proves the displayed bound for
every dimension and every nonempty compact complex matrix family of
joint spectral radius one, with

\[
\Theta_1=1,\qquad
\Theta_d=d\left(\frac{2ed^2}{d-1}\right)^{d-1}\quad(d\ge2).
\]

This constant is independent of the family and its cardinality. The
proof covers every switching word and every positive length, without
irreducibility or an exact extremal norm. The growth exponent \(d-1\) is
sharp in general; the displayed constant is not asserted optimal. The
result also holds for nonempty bounded real or complex families.

The complete original proof passed
\href{https://github.com/ajt60gaibb/OpenProblemsInNLA/blob/main/references/colbrook-jsr-growth-2026-09-11/verification/reviews/MF-05-MF-07-review.md}{independent
Codex-agent review}.
\href{https://github.com/ajt60gaibb/OpenProblemsInNLA/blob/main/references/colbrook-jsr-growth-2026-09-11/manuscripts/uniform_growth_and_holder.tex}{Authored
TeX} ·
\href{https://github.com/ajt60gaibb/OpenProblemsInNLA/blob/main/references/colbrook-jsr-growth-2026-09-11/README.md}{Submission,
authorship and verification record}. The proof was developed with AI
assistance; no external human peer review or formal verification is
claimed. The original statement and prior evidence below are retained,
and the ratings above are historical. This entry no longer contributes
to the open count.

\subsection{Context and notation}\label{context-and-notation}

Let \(\mathcal H_d\) denote the nonempty compact subsets of
\(\mathbb C^{d\times d}\).

The joint spectral radius of a nonempty compact set
\(\mathcal M\subset\mathbb C^{d\times d}\) is

\[
\widehat\rho(\mathcal M)=\lim_{k\to\infty}
\max_{A_1,\ldots,A_k\in\mathcal M}\|A_k\cdots A_1\|_2^{1/k}.
\]

This definition also applies to finite real matrix sets. The ordinary
spectral radius of one matrix is written \(\rho(A)\).

\subsection{Problem statement}\label{problem-statement}

Does each \(d\ge1\) admit \(\Theta_d>0\) such that every
\(\mathcal M\in\mathcal H_d\) with \(\widehat\rho(\mathcal M)=1\)
satisfies

\[
\|A_k\cdots A_1\|_2\le\Theta_d(Lk)^{d-1},\qquad
L=\max_{A\in\mathcal M}\|A\|_2,
\]

for all \(k\ge1\) and all \(A_1,\ldots,A_k\in\mathcal M\)?

\subsection{Reference and status
evidence}\label{reference-and-status-evidence}

Epperlein and Wirth, \href{https://arxiv.org/html/2311.18633v2}{The
joint spectral radius is pointwise Hölder continuous}, §2, Conjecture 3
(L3). Lemma 27 proves dimension two. The constant above must be
independent of the family.

\subsection{\texorpdfstring{Common follow-up screen for
\href{https://github.com/ajt60gaibb/OpenProblemsInNLA/blob/main/matrix-functions-and-stability/MF-05/README.md}{MF-05}--MF-07}{Common follow-up screen for MF-05--MF-07}}\label{common-follow-up-screen-for-mf-05mf-07}

Searches combining ``joint spectral radius'' with ``local Hölder'',
``Lipschitz lower'', ``trajectory bounds'', the authors' names, and
2025/2026 found no later resolution. These searches supplement the
explicit 2025 conjectures; they do not establish exhaustiveness. The
three entries are separately named assertions in the source, not a count
of dimensional cases.

\subsection{Audit --- 2026-09-10}\label{audit-2026-09-10}

Rechecked \href{https://arxiv.org/html/2311.18633v2}{Conjecture 3 (L3)
and Lemma 27}: dimension two is proved, but higher-dimensional uniform
constants remain conjectural. Searches for trajectory bounds and
subsequent Epperlein--Wirth work found no resolution. Family-dependent
polynomial estimates do not establish the displayed uniform statement.
\end{document}
